\documentclass[a4paper,twoside]{article}

\usepackage{morewrites}

\usepackage[svgnames,dvipsnames,x11names]{xcolor}
\usepackage{graphicx}
\usepackage[english]{babel}
\usepackage[colorlinks=true, linkcolor=NavyBlue, urlcolor=NavyBlue, citecolor=NavyBlue]{hyperref}
\usepackage[a4paper]{geometry}
\usepackage{ts-typesetting}
\usepackage{ts-fonts}
\usepackage{float}
\usepackage{xcolor}
\usepackage{epigraph}
\usepackage{marginnote}
\usepackage{multicol}
\usepackage{progressbar}
\usepackage{graphicx}
\usepackage{tcolorbox}
\usepackage{fvextra}
\usepackage{amsmath}
\usepackage{seqsplit}
\renewcommand*{\marginfont}{
\footnotesize
\sloppy
\emergencystretch=1em
\hyphenpenalty=50
\exhyphenpenalty=50
}
\newlength{\tsmarginparavail}
\newlength{\tsmarginparalt}
\newlength{\tsmarginparbuf}
\setlength{\tsmarginparbuf}{6mm}
\makeatletter
\AtBeginDocument{
\setlength{\tsmarginparavail}{\dimexpr\paperwidth-\textwidth-1in-\oddsidemargin-\marginparsep-\tsmarginparbuf\relax}
\ifdim\tsmarginparavail<\z@\setlength{\tsmarginparavail}{\z@}\fi
\if@twoside
\setlength{\tsmarginparalt}{\dimexpr1in+\evensidemargin-\marginparsep-\tsmarginparbuf\relax}
\ifdim\tsmarginparalt<\z@\setlength{\tsmarginparalt}{\z@}\fi
\ifdim\tsmarginparalt<\tsmarginparavail
\setlength{\tsmarginparavail}{\tsmarginparalt}
\fi
\fi
\ifdim\tsmarginparavail<\marginparwidth
\setlength{\marginparwidth}{\tsmarginparavail}
\fi
}
\makeatother
\usepackage{ts-code}

\usepackage{lastpage}
\usepackage{refcount}
\makeatletter
\def\ps@plain{
\let\@oddhead\@empty
\let\@evenhead\@empty
\def\@oddfoot{
\hfil
\ifnum\getpagerefnumber{LastPage}>1\relax
\thepage
\fi
\hfil}
\let\@evenfoot\@oddfoot
}
\makeatother

\title{Colorful Squares}
\author{}\date{}
\begin{document}
\maketitle
This example shows off a custom poster-ish template with four slots wired through a local template. Front matter steers colors, slot routing, and layout---no \LaTeX{} tweaks needed.
\begin{figure}[H]
\centering
\includegraphics[width=0.6\linewidth]{snippet-<HASH>.pdf}
\caption[Demo]{Demo}
\end{figure}
The YAML front matter picks the locally defined template (\tscodeinline{.}), sets the palette, and feeds each slot. Colors live under \tscodeinline{colors}, and slot content under \tscodeinline{slots}.

The manifest defines defaults, available attributes, and where they get injected.
\begin{tsdiv}{tab}[title=colorful.md]
\begin{tscode}[lang=md, engine=pygments]
\begin{Verbatim}[commandchars=\\\{\},breaklines, breakanywhere, commandchars=\\\{\}]
[include: examples/colorful/colorful.md not found]
\end{Verbatim}
\end{tscode}
\end{tsdiv}
\begin{tsdiv}{tab}[title=manifest.toml]
\begin{tscode}[lang=toml, engine=pygments]
\begin{Verbatim}[commandchars=\\\{\},breaklines, breakanywhere, commandchars=\\\{\}]
\PY{k}{[}\PY{k}{include}\PY{err}{:}\PY{+w}{ }\PY{k}{examples}\PY{err}{/}\PY{k}{colorful}\PY{err}{/}\PY{k}{manifest}\PY{k}{.}\PY{k}{toml}\PY{+w}{ }\PY{k}{not}\PY{+w}{ }\PY{k}{found}\PY{k}{]}
\end{Verbatim}
\end{tscode}
\end{tsdiv}
\begin{tsdiv}{tab}[title=template.tex]
\begin{tscode}[lang=tex, engine=pygments]
\begin{Verbatim}[commandchars=\\\{\},breaklines, breakanywhere, commandchars=\\\{\}]
[include: examples/colorful/template.tex not found]
\end{Verbatim}
\end{tscode}
\end{tsdiv}
Build it with:
\begin{tscode}[lang=text, stretch=0.5, engine=pygments]
\begin{Verbatim}[commandchars=\\\{\},breaklines, breakanywhere, commandchars=\\\{\}]
\PYZdl{} ls
colorful.md  manifest.toml  template.tex
\PYZdl{} texsmith colorful.md \PYZhy{}t. \PYZhy{}\PYZhy{}build
┌───────────────┬─────────────────────────────────────┐
│ Artifact      │ Location                            │
├───────────────┼─────────────────────────────────────┤
│ PDF           │ <STEM>.pdf                        │
└───────────────┴─────────────────────────────────────┘
\end{Verbatim}
\end{tscode}
\end{document}
